\documentclass[12 pt, a4paper]{article}% тип документа, размер шрифта
\usepackage{cmap}	
\usepackage{hyperref}
\hypersetup{
    colorlinks=true,
    linkcolor=blue,
    urlcolor=blue,
}
\usepackage{mathtext}
\usepackage[T2A]{fontenc}%поддержка кириллицы в ЛаТеХ
\usepackage[utf8]{inputenc}%кодировка
\usepackage[english,russian]{babel}
\usepackage{indentfirst}

\usepackage{amsmath,amsfonts,amssymb,amsthm,mathtools} % AMS
\usepackage{amsmath}%удобная вёрстка многострочных формул, масштабирующийся текст в формулах, формулы в рамках и др.
\usepackage{amsfonts}%поддержка ажурного и готического шрифтов — например, для записи символа {\displaystyle \mathbb {R} } \mathbb {R} 
\usepackage{amssymb}%amsfonts + несколько сотен дополнительных математических символов
\frenchspacing%запрет длинного пробела после точки
\usepackage{setspace}%возможность установки межстрочного интервала
\usepackage{indentfirst}%пакет позволяет делать в первом абзаце после заголовка абзацный отступ
\onehalfspacing%установка полуторного интервала по умолчанию
\usepackage{graphicx}%подключение рисунков
\graphicspath{{images/}}%путь ко всем рисункам
\usepackage{caption}
\usepackage{float}%плавающие картинки
\usepackage{tikz} % это для чудо-миллиметровки
\usepackage{pgfplots}%для построения графиков
\pgfplotsset{compat=newest, y label style={rotate=-90},  width=10 cm}%версия пакета построения графиков, ширина графиков
\usepackage{pgfplotstable}%простое рисование табличек
\usepackage{lastpage}%пакет нумерации страниц
\usepackage{comment}%возможность вставлять большие комменты
\usepackage{python}%питончик
\usepackage{float}
\usepackage{colortbl} % цвета в таблицах
%%%%% ПОЛЯ
\setlength\parindent{0pt} 
\usepackage[top = 2 cm, bottom = 2 cm, left = 1.5 cm, right = 1.5 cm]{geometry}
\setlength\parindent{0pt}
%%%%% КОЛОНТИТУЛЫ
\usepackage[shortlabels]{enumitem}

\usepackage{fancybox, fancyhdr}
\pagestyle{fancy} 
\fancyhead[L]{Легких Никита Б05-301}
\fancyhead[C]{2.2.1}
% ЛЮБАЯ ДОПОЛНИТЕЛЬНАЯ ИНФОРМАЦИЯ
\fancyfoot[R]{ФПМИ МФТИ}
\renewcommand{\footrulewidth}{0.3 mm}
 

\usepackage{graphicx}
\graphicspath{{pictures/}}
\DeclareGraphicsExtensions{.pdf,.png,.jpg}
\usepackage{wrapfig}
%%% Заголовок
%%% Новые команды
\newcommand{\z}[2]{{{\vspace{0.6cm} \textbf{{Задача {#1}}} \textit{{#2}\\ }}} }
\newcommand{\otv}{{\vspace{0.3cm} \textbf{Ответ: }}}
\newcommand{\uk}{\underline{\textit{Указание.}} }
\newcommand{\prim}{\underline{\textit{Примечание.}} }
\newcommand{\sol}[1]{{{\vspace{0.3cm} \textbf{{Задача {#1}} }\\ }}}



\begin{document} % конец преамбулы, начало документа

\begin{center}
\section*{Работа № 2.2.1}
    \section*{Исследование взаимной диффузии газов}
\end{center}
\textbf{Цель работы:} 1) регистрация зависимости концентрации гелия в воздухе от времени с помощью датчиков теплопроводности при разных начальных давлениях смеси газов; 2) определение коэффициента диффузии по результатам измерений.
 \\ \\
\textbf{В работе используются:} измерительная установка; форвакуумный насос; балон с газом 	(гелий); манометр; источник питания; магазин сопротивлений; гальванометр; секундомер.


\section*{Теоретическая справка}

\textit{Диффузией} называют самопроизвольное взаимное проникновение веществ друг в друга, происходящее вследствие хаотичного теплового движения молекул. При перемешивании молекул разного сорта говорят о \textit{взаимной} (или \textit{концентрационной}) \textit{диффузии}. \\

Нами будет рассматриваться именно диффузия \textit{бинарной смеси} состоящей из $He$ и обычного воздуха которая подчиняется \textit{закону Фика}: 
\begin{equation} \label{law_of_Fick}
    j = - D \cdot \nabla n,
\end{equation} 
где,  j - \textit{плотность потока компонентов} (количество частиц, пересекающих единичную площадку в единицу времени), D — \textit{коэффициент взаимной диффузии компонентов}. Знак «минус» отражает тот факт, что диффузия идёт в направлении выравнивания концентраций. \\

В нашем эксперименте будут рассматриваться такие процессы что $P$ и $T$ можно будет считать постоянными, поэтому переписав основное уравнение МКТ получаем: 
\begin{equation}\label{2}
    P = (n_{He} + n_в) k T = const \;\;\Rightarrow
\end{equation}
\begin{equation}\label{3}
    \Delta n_{He} = - \Delta n_в,
\end{equation}
где $n_{He}$ и $n_в$ — концентрации (объёмные плотности) диффундирующих газов. Исходя из \ref{3} в \ref{law_of_Fick} нам достаточно писать: \\
\begin{equation} \label{4}
    j_{He} = - D \cdot \nabla n_{He},
\end{equation} 
нами будет рассматриваться диффузия в тонкой трубке (подробности далее), поэтому в \ref{4} мы может ограничиться лишь: 
\begin{equation} \label{5}
    j_{He} = - D \cdot \frac{\partial n_{He}}{\partial x},
\end{equation}\newpage

Проведем теоретическую оценку для коэффициента диффузии. В работе концентрация гелия, как правило, мала $(n_{He} \ll n_в)$. Кроме того, атомы гелия существенно легче молекул, составляющих воздух $(\mu_{He} \ll \mu_{N_2} , \mu_{O_2})$, значит и их средняя тепловая скорость велика по сравнению с остальными частицами. Поэтому перемешивание газов в работе можно приближенно описывать как диффузию примеси лёгких частиц He на практически стационарном фоне воздуха (плюс это хорошо будет соотноситься на формулах вида \ref{5}). Коэффициент диффузии в таком приближении равен
\begin{equation}\label{6}
    D = \frac{1}{3} \lambda \overline{v},
\end{equation}
где $\overline{v} = \sqrt{\frac{8RT}{\pi \mu}}$ - средняя тепловая скорость частиц примеси, $\lambda = \frac{1}{n_0 \sigma}$. Более подробное рассмотрение показывает, что для бинарной смеси формула (\ref{6}) сохраняется, если 1) под $\lambda$ понимать величину $\lambda = \frac{1}{n_\Sigma \sigma}$, где $n_\Sigma = n_{He} + n_в = \frac{P}{kT}$ 2) $\overline{v} = \sqrt{\frac{8RT}{\pi \widetilde{\mu}}}$, где $\widetilde{\mu} = \frac{\mu_{He}\cdot\mu_{возд}}{\mu_{He}+\mu_{возд}}$  Таким образом, теория предсказывает, что коэффициент диффузии бинарной смеси \textit{обратно пропорционален давлению} в системе, и не зависит от пропорций компонентов, что мы проверим на практике.

\begin{figure}[h]
    \begin{minipage}{0.7\textwidth}
    Для исследования взаимной диффузии газов и измерения коэффициента взаимной диффузии $D$ будем использовать два сосуда объёмами $V_1$ и $V_2$ ($V_1 \approx V_2 \equiv V$), соединенные трубкой длины $L$ и сечения $S$ (\hyperref[fig:1]{рис. 1}). Предполагается, что сосуды заполнены смесью двух газов при одинаковом давлении и температуре (далее покажем как это будет осуществляться), но с различной концентрацией компонентов. Вследствие взаимной диффузии, проходящей в соединительной трубке, концентрации компонентов в сосудах с течением времени выравниваются. \\
    В общем случае концентрации компонентов $n(t,x)$ зависят от как от координаты, так и времени. Задача упрощается, если объём соединительной трубки мал по сравнению с объёмами сосудов - тогда концентрации газов $n_1(t)$ и $n_2(t)$ внутри каждого сосуда можно считать постоянными по всему объёму сосуда, и принять, что процесс выравнивания концентраций происходит благодаря диффузии в трубке.
    \end{minipage}
    \hfill
    \begin{minipage}{0.2\textwidth}
        \includegraphics[width=1\textwidth]{2.2.1/картинки/схема1.png}
        \caption*{рис. 1}
        \label{fig:1}
    \end{minipage}
\end{figure}
Для описания нашей системы решим следующую подзадачу о диффузии в соединительной трубке. Предположим сперва, что концентрации примеси (гелия) на её торцах поддерживаются постоянными и равными $n_1$ и $n_2$ соответственно. Тогда через некоторое время (оценку этого времени см. ниже ф-лу (\ref{10})) в трубке установится стационарный поток частиц, одинаковый в каждом сечении трубки (в противном случае, если бы поток зависел от $x$, частицы бы накапливались в трубке, и процесс перестал бы быть стационарным). Применяем формулу \ref{5} и получаем: 
\begin{equation}\label{7}
    j = - D \frac{\partial n}{\partial x} = const.
\end{equation}
Следовательно, распределение концентрации в трубке $n(x)$ — линейная функция:
\begin{equation*}
    n(x) = \frac{\Delta n}{L} x
\end{equation*}
и плотность потока частиц всюду постоянна и равна
\begin{equation}\label{8}
    j = - D \frac{\Delta n}{L} = - D \frac{n_2 - n_1}{L}
\end{equation} \\
В итоге мы получили зависимость $n(x)$, теперь нужно найти зависимость $n(t)$. Предположим что процесс происходит достаточно медленно чтобы считать его \textit{квазистационарным}, то есть в любой момент времени мы сможем считать что нас процесс устоявшийся и мы сможем применять закон Фика. Кроме того, будем считать, что в пределах каждого сосуда частицы распределены равномерно, так что концентрации примеси вблизи трубки и в остальных частях сосуда отличаются мало. Тогда полное число частиц примеси в сосудах равно соответственно $N_1 = n_1 V$ и $N_2 = n_2 V$ . Произведение плотности потока (\ref{8}) на площадь сечения трубки S даёт количество частиц, пересекающих в единицу времени любое поперечное сечение трубки. Поэтому 
\begin{equation*}
    \frac{dN_1}{dt} = jS,\;\; \frac{dN_2}{dt} = - jS, 
\end{equation*}
с помощью полученного и (\ref{8}) выразим $\frac{d(\Delta n)}{dt}$
\begin{equation*}
        \frac{d(\Delta n)}{dt} = - \frac{\Delta n}{\tau},
\end{equation*}
где $\tau = \frac{1}{D} \frac{VL}{2S}$, интегрируя полученное выражение получаем
\begin{equation}\label{9}
    \Delta n = \Delta n_0 \cdot e^{-t/\tau},
\end{equation}
где $\Delta n_0$ — разность концентраций примеси в сосудах в начальный момент времени. Видно, что величина $\tau$ есть \textit{характерное время} выравнивания концентраций между сосудами. Оно определяется геометрическими размерами установки и коэффициентом диффузии.

Стоит поговорить о границах применимости квазистационарного приближения, для этого обратимся к \textit{закону Эйнштейна–Смолуховского}, из которого следует нужное условие: 
\begin{equation}\label{10}
    \tau \gg \tau_{дифф} = l^2 / 2D 
\end{equation}
Кроме того, если сосуды расположены вертикально, может возникнуть вопрос о влиянии силы тяжести на диффузию. Влиянием гравитации можно пренебречь, если перепад потенциальной энергии в сосуде много меньше энергии теплового движения частиц $mgh \ll k T$ . 

\section*{Установка}
Для измерения разности концентраций в установке применяются датчики теплопроводности. При этом используется тот факт, что теплопроводность $\kappa$ смеси зависит от её состава. В общем случае зависимость $\kappa(n)$ довольно сложна, однако при малой разности $\Delta n$ концентраций в сосудах можно ожидать, что разность теплопроводностей будет изменяться прямо пропорционально $\Delta n$:
\begin{equation}\label{11}
    \Delta \kappa = \kappa(n_2) - \kappa(n_1) \approx const \cdot \Delta n
\end{equation}
Эксперименты показывают, что если доля примеси гелия составляет менее 15\%, отклонение от линейной зависимости не превышает 0,5\%, что для наших целей вполне достаточно. Устройство этих датчиков схоже с установкой из 2.2.3. \\

Для измерения сопротивлений используется мостовая схема, позволяющая определять разность показаний датчиков с высокой точностью. Мост балансируется при заполнении сосудов (и датчиков) одной и той же смесью. При заполнении сосудов смесями различного состава возникает «разбаланc» моста. При незначительном различии в составах смесей показания вольтметра, подсоединённого к диагонали моста, будут пропорциональны разности концентраций примеси: $U \sim \Delta \kappa \sim \Delta n$ . В процессе диффузии разность концентраций убывает по закону (\ref{9}), и значит по тому же закону изменяется напряжение:
\begin{equation}\label{12}
    \Delta U = \Delta U_0 \cdot e^{-t/\tau}
\end{equation}
где $U_0$ — показание гальванометра в начальный момент времени. Измеряя экспериментально зависимость $U(t)$, можно получить характерное время процесса $\tau$, с помощью которого можно рассчитать значение $D$
\begin{equation}\label{13}
     D = \frac{1}{\tau} \frac{VL}{2S}
\end{equation}
\begin{figure}[h]
\begin{center}
    \begin{minipage}{0.35\textwidth}
        Схема измерительной части установки приведена на \hyperref[fig:2]{рис. 2}. Она соединена с системой откачки и напуска воздуха и гелия. Для откачки используется форвакуумный насос. Конструкция откачки и накачки изображена на \hyperref[fig:3]{рис. 3}. Здесь воздух накачивается через кран $К_5$, гелий же накачивается с помощью двух кранов $К_6$, $К_7$, где $К_6$ - используется как дозатор, а $К_7$ как закрывающий ключ. К соединительным трубкам подключен манометр М, измеряющий разность давлений между соединительными трубками и атмосферой, и позволяющий измерять давления в разных частях системы (в зависимости от положения кранов).
    \end{minipage}
    \begin{minipage}{0.3\textwidth}
        \includegraphics[width=1\textwidth]{2.2.1/картинки/схема2.png}
        \caption*{рис. 2}
        \label{fig:2}
    \end{minipage}
    \begin{minipage}{0.3\textwidth}
        \includegraphics[width=1\textwidth]{2.2.1/картинки/схема3.png}
        \caption*{рис. 3}
        \label{fig:3}
    \end{minipage}
\end{center}
\end{figure}\newpage
Диффузия осуществляется через тонкую короткую трубку, соединяющую сосуды, оснащённую краном $К_3$. 
\begin{figure}[h]
\begin{center}
    \begin{minipage}{0.55\textwidth}
        Значения разности напряжений измеряется с помощью моста с \hyperref[fig:4]{рис. 4}. Здесь датчики теплопроводности $Д_1$ и $Д_2$, расположенные в сосудах $V_1$ и $V_2$ соответственно, включены в мостовую электрическую схему согласно \hyperref[fig:4]{рис. 4}. В одну из диагоналей моста включён высокочувствительный вольтметр (гальванометр) Г, к другой подключается источник небольшого постоянного напряжения ($\mathcal{E} = 2.1\;В$). Сопротивления проволок датчиков составляют одно из плеч моста. Второе плечо составляют переменные сопротивления $R_1$, $R_2$ и $R$, служащие для установки показаний вольтметра Г на нуль (балансировка моста).
    \end{minipage}
    \begin{minipage}{0.35\textwidth}
        \includegraphics[width=1\textwidth]{2.2.1/картинки/схема4.png}
        \caption*{рис. 4}
        \label{fig:4}
    \end{minipage}
\end{center}
\end{figure}

\section*{Задание}
\textbf{1.}  
\begin{enumerate}
    \item \hyperref[fig:2]{См. тут}
    \item В нашем случае $p_{атм} = 758.5\;мм\;рт.\;ст. \equiv 102\;делений \rightarrow \gamma = 7.4\;мм\;рт.\;ст/дел.$ Объем сосудов для исследуемых газов $V = (800\pm5)\;мл$. Характеристика трубы $\frac{L}{S} = (15.0\pm0.1) \frac{1}{см}$
    \item Запустили, создали папку folder
\end{enumerate}
\textbf{2.}
Установка подготовлена к работе \\
\textbf{3 - 6.}
Измерения проведены в соответствии с порядком в \hyperlink{https://old.mipt.ru/education/chair/physics/S_II/lab/rabota221/221.pdf}{лабнике}, результаты измерений находятся \hyperlink{https://drive.google.com/drive/u/0/folders/18QGdCFIDWtZ8oKLtq0SHHksDOskPXUkK}{тут} \\
\textbf{7.} В этом пункте за значение $P_\Sigma$ взято число 40.9 торр, $P_{He} = 36.8$ торр, $P_{возд} = 4.1$ торр. Полученные измерения лежат \hyperlink{https://drive.google.com/drive/u/0/folders/18QGdCFIDWtZ8oKLtq0SHHksDOskPXUkK}{тут} \\
\section*{Обработка результатов измерений}
\textbf{8.} Нашей целью является проверка закономерности (\ref{9}), для этого воспользуемся тем что
\begin{equation}\label{14}
    (\ref{9}) \Rightarrow (\ref{12}) \Rightarrow ln\left(\frac{U}{U_0}\right) \sim - \frac{t}{\tau},
\end{equation}
поэтому будем проверять (\ref{14}), для этого построим логарифмированные графики для каждого рабочего давления и проверим их на линейность.
\begin{figure}[h]
    \begin{center}
        \begin{minipage}[h]{0.49\linewidth}
        \includegraphics[width=1\linewidth]{2.2.1/картинки/40.9.png}
        \caption*{Угол наклона: $k_1 = -(2.142\pm0.001\ 1/мс \rightarrow$ \\ $\rightarrow \tau = (467.80\pm0.18)\;с \rightarrow$\\ $\rightarrow D = (12.85\pm0.17)\;см^2/c$}
        \label{fig:40.9} %% метка рисунка для ссылки на него
        \end{minipage}
        \hfill
        \begin{minipage}[h]{0.49\linewidth}
        \includegraphics[width=1\linewidth]{2.2.1/картинки/74.4.png}
        \caption*{Угол наклона: $k_2 = -(1.22\pm0.001)\ 1/мс \rightarrow$ \\ $\rightarrow \tau = (820.74\pm0.44)\;с \rightarrow$\\ $\rightarrow D = (7.31\pm0.10)\;см^2/c$}
        \label{fig:74.4}
        \end{minipage}
    \end{center}
\end{figure} \newpage
\begin{figure}[h]
    \begin{center}
        \begin{minipage}[h]{0.49\linewidth}
        \includegraphics[width=1\linewidth]{2.2.1/картинки/104.1.png}
        \caption*{Угол наклона: $k_4 = -(0.814\pm0.001)\ 1/мс \rightarrow$ \\ $\rightarrow \tau = (1228\pm1)\;с \rightarrow$\\ $\rightarrow D = (4.88\pm0.07)\;см^2/c$}
        \label{fig:104.1} %% метка рисунка для ссылки на него
        \end{minipage}
        \hfill
        \begin{minipage}[h]{0.49\linewidth}
        \includegraphics[width=1\linewidth]{2.2.1/картинки/148.7.png}
        \caption*{Угол наклона: $k_4 = -(0.605\pm0.001)\ 1/мс \rightarrow$ \\ $\rightarrow \tau = (1652\pm2)\;с \rightarrow$\\ $\rightarrow D = (3.63\pm0.05)\;см^2/c$}
        \label{fig:148.7}
        \end{minipage}
    \end{center}
\end{figure} \newpage
\begin{figure}[h]
    \centering
    \includegraphics[width=0.5\linewidth]{2.2.1/картинки/185.9.png}
    \caption*{Угол наклона: $k_5 = -(0.479\pm0.001)\ 1/мс \rightarrow$ \\ $\rightarrow \tau = (2087\pm2)\;с \rightarrow$\\ $\rightarrow D = (2.86\pm0.04)\;см^2/c$}
    \label{fig:185.9}
\end{figure} \\

\textbf{9.} Объеденим значения из прошлого пункта в таблицу и построим требуемый график: 
\begin{center}
\begin{tabular}{|
>{\columncolor[HTML]{D9D2E9}}c |
>{\columncolor[HTML]{D9D2E9}}c |
>{\columncolor[HTML]{D9D2E9}}c |
>{\columncolor[HTML]{D9D2E9}}c |}
\hline
$P$, торр & $1/P$,  торр$^{-1} \cdot 10^{-3}$ & $D$,  см$^2$/c & $\Delta D, \;см^2/c$ \\ \hline
40.9      & 24.45                           & 18.85          & 0.17                 \\ \hline
74.4      & 13.44                           & 7.31           & 0.10                 \\ \hline
104.1     & 9.61                            & 4.88           & 0.07                 \\ \hline
148.7     & 6.72                            & 3.63           & 0.05                 \\ \hline
185.9     & 5.38                            & 2.86           & 0.04              
  \\ \hline
\end{tabular}
\end{center}
\begin{figure}[h]
    \centering
    \includegraphics[width=0.7\linewidth]{2.2.1/картинки/D(frac(1)(P)).png}
    \label{fig:D(frac(1)(P))}
\end{figure} \newpage

\begin{equation*}
    Угол\;наклона:\;\;k = (0.85\pm0.05)\;см^2\cdotторр/с,
\end{equation*}
\begin{equation*}
    Свободный\;коефициент:\;\;b = (-2.6\pm0.6)\;см^2/c\;\;\rightarrow
\end{equation*}
\begin{equation*}
    \rightarrow\;\;D_{атм} = -(1.5\pm0.7)\;см^2/c.
\end{equation*}
Мы получили явно неадекватное значение, как минимум из-за знака. Визуальный осмотр графика дает предположение что ошибка таится в самой первой точки, попробуем построить график без нее и посмотрим на результат:
\begin{figure}[h]
    \centering
    \includegraphics[width=0.7\linewidth]{2.2.1/картинки/D(frac(1)(P))_fixed.png}
    \label{fig:D(frac(1)(P))_fixed}
\end{figure}
\begin{equation*}
    Угол\;наклона:\;\;k = (0.54\pm0.03)\;см^2\cdotторр/с,
\end{equation*}
\begin{equation*}
    Свободный\;коефициент:\;\;b = (-0.10\pm0.24)\;см^2/c\;\;\Rightarrow\;свободным\;коефцициентом\; можно\;пренебречь,\;причем\;это\;сходится\;с\;теорией
\end{equation*}
\begin{equation*}
    \;причем\;это\;сходится\;с\;теорией\;\;\rightarrow
\end{equation*}
\begin{equation*}\label{D_atm}
    \rightarrow\;\;D_{атм} = (0.61\pm0.03)\;см^2/c.
\end{equation*}
Табличное значение равно $D_{табл.\;атм} = 0.66\;см^2/c$, значит мы получили значение очень близкое к правде, в свою очередь $\left( 1 - \frac{D_{атм}}{D_{табл.\;атм}}\right) = 8 \%$. Остановимся на этом результате потому что он не сильно отстоит от реальности.
\newpage
\textbf{10.} Строим нужный график:
\begin{figure}[h]
    \centering
    \includegraphics[width=0.5\linewidth]{2.2.1/картинки/77.7.png}
    \caption*{Угол наклона: $k_7 = -(1.28\pm0.06)\ 1/мс \rightarrow$ \\ $\rightarrow \tau = (780\pm4)\;с \rightarrow$\\ $\rightarrow D = (7.69\pm0.14)\;см^2/c$}
    \label{fig:77.7}
\end{figure} \\
Исходя из данных графиков: $D_{He-возд} = (12.85\pm0.17)\;см^2/c,\;D_{возд-He} = (7.69\pm0.14)\;см^2/c$. \\
Считая что температура и давление в обоих случаях одинаковые должно выполняться следующее соотношение: 
\begin{equation}\label{15}
    \frac{D_{He-возд}}{D_{возд-He}} = \sqrt{\frac{\mu_{возд}}{\mu_{He}}}
\end{equation}
В нашем случае получается что: 
\begin{equation*}
    \frac{D_{He-возд}}{D_{возд-He}} = 1.7 \approx 2.7 = \sqrt{\frac{\mu_{возд}}{\mu_{He}}}
\end{equation*}
\begin{equation*}
    \left( 1 - \frac{\frac{D_{He-возд}}{D_{возд-He}}}{\sqrt{\frac{\mu_{возд}}{\mu_{He}}}}\right) \cdot 100\% = \left( 1 - \frac{1.7}{2.7}\right) \cdot 100\% = 35\%
\end{equation*}
Погрешность получается огромной, но это можно связать с тем что в прошлом пункте измерение при 40.9 торр тоже было "лишним". \\
\textbf{11.} Теперь по полученным данным оценим длину свободного пробега и размер молекулы. $T_{атм} = 23^oC = T_{He}$. Поскольку 
$D_{} = \frac{1}{3} \lambda \overline{v_{He}},\;$ где $\overline{v_{He}} = \sqrt{\frac{8RT_{He}}{\pi \mu}},$
то длина свободного пробега $\lambda_{He} = 3D_{He}\sqrt{\dfrac{\pi \mu}{8RT_{He}}} \approx (163\pm2)$ нм, что сходится с табличными значением $\lambda_{атм} = 175$ нм.


Эффективное сечение  $\sigma_{He} = \dfrac{kT_{атм}}{\sqrt{\lambda P_{атм}}} \approx 23 \cdot 10^{-2} $ $нм^2$.

\section*{Вывод}
Проведя регистрацию зависимости концентрации гелия в воздухе от времени с помощью датчиков теплопроводности при разных начальных давлениях смеси газов мы экспериментально подтвердили зависимость (\ref{9}) с помощью его связи с (\ref{12}) и нашли значения коэффициентов диффузии для различных рабочих давлениях. Также была подтверждена линейная зависимость $D(\frac{1}{P})$ и было получено значение $D_{атм}$, значение которого схоже с табличным (\hyperref[D_atm]{подробнее тут}).

Также мы нашли значения $\lambda_{He} = 175$ нм, $\left( 1 - \frac{163}{175}\right) \cdot 100\% \approx 7\%$  и значение $\sigma_{He} = 23 \cdot 10^{-2} $ $нм^2$

\end{document}